%% Supplementary material for the PhysDDG manuscript.
%% Generated from supplementary.md (the markdown twin); keep the two in step.
%% Compile with: tools/tectonic -X compile supplementary.tex
\documentclass[11pt,a4paper]{article}
\usepackage[utf8]{inputenc}
\usepackage[T1]{fontenc}
\usepackage[margin=2.2cm]{geometry}
\usepackage{booktabs}
\usepackage{array}
\usepackage{amsmath}
\usepackage{longtable}
\usepackage[hidelinks]{hyperref}
\usepackage{parskip}
\renewcommand{\arraystretch}{1.1}
\setcounter{secnumdepth}{0}
\title{Supplementary: methods that were tried and did not work\\[2mm]
\large PhysDDG: a GPU physics engine with a linear read-out for protein stability changes}
\author{<AUTHOR PLACEHOLDER>}
\date{}
\begin{document}
\maketitle
\sloppy


Supplement to "PhysDDG: a GPU physics engine with a linear read-out for protein stability changes".

Every number here comes from the project's own result tables (\texttt{results/ladder\_physics/},
\texttt{results/af\_mode/}, \texttt{results/enm\_flexibility/}, \texttt{results/diff\_params/}, \texttt{results/evo\_extension/}) and from the
notes that describe them. Confidence intervals are 95\% percentile intervals from a bootstrap that resamples
homology clusters. Comparisons are paired on identical rows. "Val" means the Megascale validation split and
"CV" means grouped five-fold cross-validation over homology clusters on the training split; those are the
numbers that no test set influenced.

Shorthand for the models: \textbf{base Stage 1} and \textbf{base Stage 2} are the fitted engine models without the two
late terms; \textbf{default} is the shipped Stage-2 model; \textbf{joint} is the variant fitted on Megascale and FireProt
together. The Megascale primary subset is the 21 natural test proteins (22,159 variants).



\section{S1. Reweighting the 22 FoldX energy terms}

\textbf{Hypothesis and source.} FoldX publishes its 22 per-term energies for every prediction
(Guerois 2002; Schymkowitz 2005). If its weights, fitted to about a thousand measurements in 2002, are the
limitation, refitting them on 215,731 Megascale rows should help.

\textbf{What was done.} Free linear weights on the 22 FoldX terms, fitted under protein-grouped five-fold
cross-validation on Megascale-train, with and without log compression of the clash terms.

\textbf{Numbers.} Pooled Spearman on Megascale-test: stock FoldX 0.569; refitted 0.541; refitted with
log-compressed clash terms 0.567 (per-protein 0.633 / 0.604 / 0.628). Seven of the 22 terms are identically
zero on monomeric single mutants (loop entropies, water bridge, disulfide, electrostatic kon, partial covalent,
complex entropy) and three more have $|$rho$|$ <= 0.02 (cis-bond, ionisation, helix dipole).

\textbf{Why we think it failed.} The missing accuracy is in the terms and in the sampling, not in the weights.
Refitting more parameters on a different assay than the one that FoldX was tuned on loses more than it gains.

\textbf{Decision.} Do not build on FoldX terms; build the engine, and keep the term count near 20.



\section{S2. Hard cap versus soft repulsion, and site-dependent caps}

\textbf{Hypothesis and source.} A fixed backbone over-penalises small-to-large substitutions, and the repulsion term
is both the heavy tail and the antisymmetry defect. The design critique proposed a site-dependent cap
\texttt{cap = c0 + c1 * flex(site)} (decision D9), with flexibility from the smoothed contact number.

\textbf{What was done.} Three cap settings (capA \texttt{c0=1}; capB \texttt{c0=3}, the original; capC \texttt{c0=1, c1=2} rising with
burial), each with an additional log-softened repulsion column, with all features recomputed. Selection on
Megascale-val plus grouped CV only.

\textbf{Numbers.} Val natural per-protein rho: capA hard 0.7276, capB hard 0.7265, capC hard 0.7291 (selected),
soft variants 0.7190--0.7201. CV: 0.6672--0.6715. The spread over all six settings is 0.010 on validation and
0.004 in CV, inside the CV standard error of 0.0064. On test, refitting with the winner moves the Stage-1 model
from 0.677 to 0.678. The soft column halves the heavy tail (99th percentile of the repulsion feature 123 to 53
kcal/mol) at a cost of about 0.007 validation rho. On the antisymmetry diagnostic the unit-weight repulsion bias
is +7.5 kcal/mol (capB), +4.2 (capA) and +3.2 (soft column); the correlation between direct and negated reverse
stays at 0.29--0.35.

\textbf{Why we think it failed.} The cap controls the size of the tail, not its cause. Each structure is packed
around the residue that it contains, so both directions of a substitution clash, and the fixed backbone is the cause.

\textbf{Decision.} Keep capC on physical grounds and state that the choice does not matter. The signed-log transform
of the repulsion feature stays, selected on validation.



\section{S3. Single-structure symmetrisation and antisymmetry fixes}

\textbf{Hypothesis and source.} A physics model should obtain antisymmetry for free (Pucci 2018). Three fixes were
tried.

\textbf{What was done and the numbers.}

\begin{center}\scriptsize
\begin{tabular}{@{}>{\raggedright\arraybackslash}p{0.409\linewidth}>{\raggedright\arraybackslash}p{0.153\linewidth}>{\raggedright\arraybackslash}p{0.062\linewidth}>{\raggedright\arraybackslash}p{0.277\linewidth}@{}}
\toprule
fix & correlation(direct, $-$reverse) & bias (kcal/mol) & cost \\
\midrule
none (default model) & 0.25 & +1.37 & --- \\
symmetric residue-pair offsets S(a,b) fitted inside the reference energies & 0.59 vs 0.587 without, i.e. no change & +0.50 & none, and no benefit \\
drop every wild-type-site-only main effect ("odd" model) & 0.57--0.59 & +0.38--0.44 & $-$0.044 Megascale per-protein rho \\
two-structure symmetrisation $\tfrac{1}{2}$[pred(direct) $-$ pred(reverse on the mutant crystal)] & $-$1 by construction & 0 & direct Pearson 0.70 to 0.60, and needs the mutant crystal structure \\
single-structure symmetrisation: score the reverse on the modelled mutant environment (\texttt{sym\_s2\_dF}) & 0.15 & +3.5 (unfitted) & Megascale-val raw rho 0.558 to 0.540 \\
\bottomrule
\end{tabular}\end{center}

Per term, every engine energy is already nearly antisymmetric across the two structures (correlation 0.90--0.98,
mean direct + reverse $\pm$0.1--0.3 kcal/mol). The exception is repulsion: mean direct + reverse +7.5 (Stage 1) /
+5.9 (Stage 2) kcal/mol, correlation 0.29.

\textbf{Why we think it failed.} The modelled mutant environment stays close to the wild-type crystal environment,
because the backbone is fixed and the neighbours move little, while the real mutant crystal is repacked. A
single structure therefore cannot deliver cross-structure antisymmetry in a fixed-backbone model.

\textbf{Decision.} Report the default model as not antisymmetric, offer the "odd" variant where antisymmetry
matters, and treat Ssym as a diagnostic only (96.5\% of its rows have a FireProt training homologue).



\section{S4. Backbone ensembles (gate G3)}

\textbf{Hypothesis and source.} Backbone give is the physically correct fix for over-penalised small-to-large
substitutions; ensemble averaging helps FoldX and is the whole advantage of Rosetta flex-ddG (Barlow 2018).

\textbf{What was done.} Native plus eight perturbed backbones per structure; each residue displaced by a smooth
Gaussian-process field (RBF length 6 \AA{}) at 0.3 \AA{} or 0.6 \AA{} RMS amplitude; torsions and rotamer candidates kept
from the native structure; aggregation by mean and by Boltzmann weight over backbones. Megascale validation
only, 17,908 natural rows, Stage-1 features, weights not refitted.

\textbf{Numbers.} Per-protein rho: native 0.7265; 0.3 \AA{} mean 0.7257; 0.3 \AA{} Boltzmann 0.7253; 0.6 \AA{} mean 0.7233;
0.6 \AA{} Boltzmann 0.7188. Buried small-to-large (1,402 rows): 0.349 native, at most 0.358 with an ensemble,
against a gate of +0.015. Cost: nine times Stage 1.

\textbf{Why we think it failed.} Random displacements do not find the relaxation that matters. The relevant motion is
a directed accommodation of the new side chain, which needs a minimiser rather than noise.

\textbf{Decision.} Dropped. A directed local relaxation remains the untested option and is the first thing we would
try next.



\section{S5. Neighbour repacking overall (gate G2)}

\textbf{Hypothesis.} Letting the neighbours move should matter wherever the mutation changes the packing.

\textbf{What was done.} Stage 2 as described in the main text, compared with Stage 1 on identical rows and, after
audit 1 objected to a test-based stratum claim, re-done on out-of-fold predictions on the training split and on
validation.

\textbf{Numbers.} Overall gain: +0.004 (CV out of fold), +0.005 (validation), +0.010 on test, against a
pre-registered gate of +0.02. On buried small-to-large substitutions (8--9\% of rows): +0.034 pooled and +0.058
per-protein out of fold, in 5 of 5 folds, and +0.068 / +0.078 on validation, which clears the +0.03 gate for
that stratum. Runtime: about 550 mutations/s/GPU against 1,900 for Stage 1. A diagnostic worth noting: the
mean-field packer returns the neighbour's native rotamer as its most probable state in only about half of the
wild-type cases, partly because the native rotamer and its library twin split the probability.

\textbf{Decision.} Gate failed overall, passed on its target stratum. Stage 2 is the default for crystal input
because it is affordable there, Stage 1 is the default for proteome-scale scans, and the honest statement is
that repacking buys accuracy in one stratum rather than everywhere.



\section{S6. Elastic-network flexibility as a dynamics term}

\textbf{Hypothesis and source.} 43\% of the residual variance is at the site level and no descriptor in the model
describes site dynamics. Gaussian and anisotropic network models give site flexibility for free
(Bahar 1997; Atilgan 2001).

\textbf{What was done.} Eleven descriptors (GNM mean-square fluctuation, ANM, centroid-node ANM, Eyal--Bahar spring
constants, perturbation-response scanning, dynamic flexibility index, hinge score, contact degree, a
pLDDT-weighted GNM), three normalisations and eight functional forms, so 264 configurations per stage, with a
pre-registered primary candidate (GNM fluctuation multiplying the engine energy) and a pre-registered bar
(+0.004 on validation and in CV with at least 4 of 5 folds up).

\textbf{Numbers.} The primary candidate is rejected: validation $-$0.005, 1 of 5 folds up at Stage 2. The best of the
264 (centroid ANM fluctuation rank scaling the whole prediction, weight $-$0.56) passes the bar: validation
+0.005 / +0.011, CV +0.012 / +0.014, and on the single test look Domainome +0.012 [+0.009, +0.023] and
Megascale natural-strict +0.015 [$-$0.000, +0.038]. But a static control, the contact-number rank in the same
form, recovers about 75\% of it (validation +0.003, CV +0.009), and adding the dynamic descriptor on top of the
static control is worth +0.002 / +0.003. Every elastic-network descriptor correlates $-$0.85 to $-$0.91 with the
contact number. On the site-level residuals the dynamic descriptors add at most 0.004 R$^2$ over the static set.
The term also costs buried small-to-large (0.487 to 0.469), from-Gly (0.645 to 0.615) and antisymmetry (0.265
to 0.162), and it calls the buried-but-tolerant SH3 positions rigid, so it makes exactly the sites that it was meant
to fix slightly worse.

\textbf{Decision.} Not adopted (D21). It is a packing-dependent scale on the prediction, not dynamics, and it is
chosen from 264 configurations; the honest version is the contact-number control.



\section{S7. Fitting the 32 physical parameters end to end}

\textbf{Hypothesis and source.} Every physical constant is a frozen literature value. Rosetta-style refits of the
energy parameters themselves usually buy accuracy (Park 2016; Alford 2017).

\textbf{What was done.} 32 multiplicative log-scale parameters (14 EEF1 solvation strengths, the EEF1 correlation
length, four radii, four well depths, four term weights inside the packing energy, the Debye length, the
ensemble temperature, the rotamer-prior weight, the repulsive-radius scale and the repulsion cap), fitted
through the differentiable Stage-1 scan with the Boltzmann average included, using variable projection: for
every parameter vector the full linear model is solved to optimality and the gradient is the envelope-theorem
gradient, checked against finite differences to better than 0.5\%. Trained on a site-balanced subsample of
Megascale-train (14,340 rows); the penalty and stopping iteration chosen by grouped CV plus validation.

\textbf{Numbers.} Training loss falls 1.9\%. Validation per-protein rho 0.7262 to 0.7260 (flat); CV natural
per-protein +0.0055 but in only 3 of 5 folds; pre-set bar (+0.01 on validation and 4 of 5 folds) not met. On the
single test look: Megascale natural +0.005 [$-$0.007, +0.016] rank, Domainome $-$0.001, FireProt +0.011
[$-$0.037, +0.048] Pearson; the only consistent signal is calibration (Megascale Pearson +0.005
[+0.002, +0.009], and positive on every set).

\textbf{What it tells us anyway.} Every fitted parameter stays within a factor of 1.6 of its literature value, none
hits its bound, and a Megascale-only and a FireProt-only fit agree at r = 0.80 over the 32 parameters. Both fits
say the same thing: polar radii $\times$0.77, polar well depth $\times$0.74, backbone-nitrogen desolvation $\times$0.62, carboxyl and
backbone-oxygen desolvation $\times$1.2--1.3. The literature physics is already close to a local optimum for this task.

\textbf{Decision.} Nothing shipped; reported as a negative result with a transferable but small calibration gain.



\section{S8. Floor-aware loss}

\textbf{Hypothesis and source.} The second critique noted that 4,010 training rows (1.8\%) sit at the per-domain
measurement floor and carry 6.6\% of the squared error; treating them as lower bounds should remove a bias.

\textbf{What was done.} Refit the default model with those rows as Tobit lower bounds, using the same interval
machinery already applied to the 11,363 excluded rows.

\textbf{Numbers.} Validation natural per-protein rho 0.7411 to 0.7412; grouped CV 0.6813 to 0.6718; no weight moves
by more than 3.8\%, against the critique's own diagnostic of 20\%.

\textbf{Decision.} Not adopted. The fit is not floor-biased, and censoring those rows discards real information.
The separate, successful use of censoring (the 11,363 rows that the assay excluded) stays in every fit.



\section{S9. Pairwise and double-mutant terms}

\textbf{Hypothesis.} Explicit coupling terms could capture non-additivity, and Megascale contains 127,476 double
mutants whose two singles are also measured, so the non-additivity epsilon = ddG(AB) $-$ ddG(A) $-$ ddG(B) is
measurable without any model.

\textbf{Numbers.} Raw, non-additivity looks enormous and distance-independent: mean epsilon $-$0.88 kcal/mol, $|$epsilon$|$
> 0.5 for 64\% of doubles, and the mean is as negative beyond 15 \AA{} ($-$0.76) as below 5 \AA{} ($-$0.78). That is the assay
floor, not coupling. Restricting to doubles that cannot be floor-limited (sum of singles < 1.5 kcal/mol,
28,123 rows) leaves mean epsilon $-$0.225 with sd 0.50 and $|$epsilon$|$ > 0.5 in 31\%, and the distance trend is not
monotone: the contact bin below 5 \AA{} has the \emph{smallest} mean ($-$0.093) and the 6.5--8 \AA{} shell the largest ($-$0.321).

\textbf{Decision.} No pairwise term. The upper bound on what one could buy is about 0.5 kcal/mol rms on doubles and
nothing identifiable on singles, which is what every benchmark scores. A double-mutant scan entry point in the
engine is specified but not built; that would be a capability argument, not an accuracy one.



\section{S10. The two site terms that failed transfer (gate G6)}

\textbf{Hypothesis.} The residual analysis found that per-site sensitivity is where the error is, and proposed two
wild-type-site descriptors: side-chain contact number at the site (\texttt{site\_scsc}) and the number of buried
side-chain hydrogen-bond partners (\texttt{site\_hbsc\_bur}).

\textbf{Numbers.} Both help on Megascale: validation 0.7329 to 0.7438 (\texttt{site\_scsc}) and CV 0.6727 to 0.6845
(\texttt{site\_hbsc\_bur}). Transfer test (weights fitted on Megascale only against FireProt only, ratio scaled by the
1.7 assay-scale factor): \texttt{site\_hbsc\_bur} 0.25 --- fails the three-fold rule, and is undetermined on FireProt
(0.039 $\pm$ 0.047). \texttt{site\_scsc} 0.35, at the edge of the rule. On the crystal benchmarks \texttt{site\_scsc} costs FireProt
0.464 to 0.439 Pearson, S669-strict 0.316 to 0.291 and S461-strict 0.549 to 0.545, while adding +0.009 on
Megascale and +0.019 on Domainome. It also pulls the burial weight down from 0.932 to 0.536, i.e. it takes over
part of the burial main effect.

\textbf{Decision.} Both dropped from the shipped model (D19). The model with \texttt{site\_scsc} is kept in the tables for
transparency. This decision used FireProt and S669 test numbers and is one of the five test-informed decisions
disclosed in S16.



\section{S11. The disulfide term, eight rows, and the catalytic-CXXC exception}

This is the clearest failure in the project and it took two audits to resolve.

\textbf{What happened.} Megascale contains no disulfide-breaking mutations, so a \texttt{ss\_broken} indicator can only be
fitted on FireProt, where it rests on about 8 training rows. It came out at +3.34 $\pm$ 1.04 kcal/mol, which is a
sane penalty for a structural disulfide and badly wrong for a catalytic one. Thioredoxin 2TRX C32A and C35A
(the redox-active CGPC motif) are measured at $-$4.17 and $-$4.00 kcal/mol --- breaking that disulfide \emph{stabilises}
the protein --- and the joint model predicted +4.31 and +3.81.

\textbf{The first wrong fix.} A CXXC exception was added that switched the term off for catalytic motifs. Audit 1
found that this rule had been derived from the two FireProt \emph{test} rows it fixes: with the exception the joint
model scores Pearson 0.479--0.484 on FireProt-test, without it 0.389--0.396. The exception was removed from the
headline.

\textbf{The second wrong fix.} The reporting script then evaluated FireProt on the rows that every predictor covers,
which is 338 of 347 --- and the two rows missing from one baseline were exactly the 2TRX pair. Audit 2 found this.
On the honest row set the joint model's FireProt Pearson is 0.388, not 0.482: below the Megascale-only models
(0.451--0.456) and below clipped FoldX (0.417). The whole FireProt advantage of the joint fits was those two
dropped rows.

\begin{center}\scriptsize
\begin{tabular}{@{}llllll@{}}
\toprule
FireProt-test row set & n & joint & default & base Stage 2 & FoldX-clip \\
\midrule
as previously reported & 338 & 0.482 & 0.464 & 0.467 & 0.423 \\
our models + FoldX & 341 & 0.388 & 0.453 & 0.458 & 0.417 \\
all rows & 347 & 0.388 & 0.451 & 0.456 & 0.417 \\
all rows, whole 2TRX protein removed & 328 & 0.462 & 0.442 & 0.445 & 0.406 \\
\bottomrule
\end{tabular}\end{center}

\textbf{The fix that stands.} The \texttt{ss\_broken} term is dropped from the joint fits by the same transfer rule that had
already removed \texttt{site\_hbsc\_bur}, decided on validation, CV and transfer only (\texttt{ss\_decision.csv}: dropping it
changes validation 0.7285 to 0.7297 and CV 0.6607 to 0.6631). FireProt is now reported on all 347 rows in every
table, and any comparison that omits rows says so.

\textbf{Lesson.} A term fitted on eight rows needs a scope, and the scope must come from training data or from the
literature. An "identical rows" rule that intersects predictor coverage can silently undo a correction.



\section{S12. AlphaFold treatments that did not ship}

\textbf{pLDDT shrinkage.} A two-parameter logistic in pLDDT multiplying the prediction gains +0.006 to +0.022
per-domain rho on Domainome, which passes the literal gate for one of the four predictors tested. The
burial-matched control --- the identical logistic in 100$\cdot$(1 $-$ RSA), fitted on the same held-in data --- gains
+0.014 to +0.073, i.e. up to three times as much. The treatment is a burial-dependent calibration that carries a
pLDDT label; on Domainome, pLDDT is largely a burial label, because low-pLDDT sites are almost all exposed.
\emph{Decision:} not shipped as an AlphaFold correction (D8). A later, different version of this term survived a
re-run of the same control inside a model that already contains burial and burial-weighted terms (there the
sign-constrained burial control is pinned at zero, and refitting it without the sign constraint gives +0.82,
which amplifies buried predictions instead of attenuating them), and it ships for AlphaFold inputs only with its weight interval including zero ($-$0.97, bootstrap
interval $-$2.52 to $-$0.23 over 200 draws, and the model-level gain +0.012 [+0.006, +0.019] on Domainome).

\textbf{pLDDT-modulated repulsion cap.} Exactly 0.000 on Domainome. A uniformly softer cap is harmful: $-$0.004 for
the fitted model and $-$0.028 for the reference physics model.

\textbf{Global repacking of all side chains.} Agreement between the crystal and the AlphaFold prediction falls
(pooled r 0.926 to 0.921; fraction of entities above 0.9: 0.709 to 0.636) and accuracy does not move
(0.573 to 0.571). Small backbone differences get amplified into different rotamer choices. This reproduces
Vangaru and Bhattacharya (2025) in our own energy: keep the AlphaFold side chains, repack only the site and its
shell.

\textbf{Backbone-only burial in the 70--90 band.} +0.004 for one predictor, $-$0.009 for another; below the gate.

\textbf{Trimming (the treatment that is kept).} The MAVISp-style rule gains 0.000 to +0.009 on Domainome and +0.014
to +0.022 on Megascale-test, at the cost of 0.9\% of sites. Every Domainome interval touches zero. It is kept
because it is free, standard and never harmful, not because it is proven. The Megascale gain is a test-set
observation (the AlphaFold parameters were chosen on abundance data and half of Domainome), so it is supporting
evidence only.

\textbf{What actually works on AlphaFold input: flagging.} Withholding variants with site pLDDT below 70, within 3
residues of a construct terminus, or in an exposed hydrophobic patch (14\% of rows) gains +0.038 per-domain rho
for the fitted engine model, +0.035 for the reference physics model and +0.020 for the descriptor baseline ---
more than every computation-changing treatment together. Withholding metal sites, by contrast, \emph{lowers} the
score: those 7\% of rows are still partly predictable from burial, and they are the one place where FoldX
collapses (0.44 to 0.14) while the physics models keep 0.22--0.39. Metal sites should be flagged and reported,
not withheld.



\section{S13. tmol as the engine}

\textbf{Hypothesis and source.} tmol (Apache-2.0) is the Rosetta energy, Dunbrack packer and minimiser in batched
PyTorch and CUDA, with no neural network. If it is fast and accurate enough, writing our own kernels is waste.

\textbf{Numbers.} Raw tmol ddG is clearly worse than FoldX on Megascale (pooled Spearman 0.463 repack-only, 0.492
with torsion minimisation, against FoldX 0.569), because 17\% of variants have a repulsion change above 20 REU
from fixed-backbone clashes and the Pearson correlation collapses to 0.18--0.21. Refitting its 12 terms plus 20
reference energies on a Megascale training subset, with all choices made by grouped CV, lifts it to pooled
0.596 / per-protein 0.626 --- the level of FoldX and of our 10-weight baseline, and below our engine (0.692
per-protein). Speed: 3.1 minutes for a 349-residue scan, against 18.7 s for our Stage 2 on 490 residues.

\textbf{Decision.} Not adopted as the engine; kept as a baseline, where it is a useful one because it correlates only
0.76 with our descriptor baseline and 0.66 with FoldX. Its bundled Dunbrack rotamer library has a licence
question we did not resolve.



\section{S14. Joint Megascale + FireProt fitting versus Megascale only}

\textbf{Hypothesis.} Fitting both assays with per-assay scale and offset should buy generality.

\textbf{What was done and the numbers.} The joint family is better on the abundance-type assays (Domainome 0.450
against 0.418; strict 0.403 against 0.391; the seven abundance assays 0.484 against 0.469) and on FireProt
(0.478 against 0.451 on all 347 rows), at a visible cost on the crystal sets (S669-strict 0.319 against 0.336;
S461-strict 0.531 against 0.552) and no measurable Megascale cost ($-$0.001 [$-$0.038, +0.016] pooled). None of
the FireProt differences is significant on any row set.

\textbf{The decision rule, fixed before looking and applied to validation, CV and transfer only:} ship the family
that wins on (i) Megascale validation natural per-protein rho, both families scored on identical held-out rows,
which is the only directly comparable number that is not a test number; then (ii) grouped CV on each model's
own training data; then (iii) transfer. Ties go to fewer fitted terms. Outcome: Megascale-only 0.7411
validation / 0.6813 CV against joint 0.7297 / 0.6648 --- \textbf{the Megascale-only family ships}, and the joint fits
ship as the multi-assay variant.

\textbf{Honest note.} An earlier version of this project shipped the joint family, on a gate that was stated
directly on Domainome and Megascale test numbers. Audit 2 flagged that, and the decision above is the
re-decision without test data. Also note that FireProt-test is not an out-of-assay test for the joint models at
all, and that 14 of its 347 rows share a single-linkage family with a FireProt training protein.



\section{S15. Domainome as a target, and the reference energies that are wrong there}

\textbf{What we found.} Domainome measures in-cell abundance of isolated domains, not folding free energy. Per-mutant
amino-acid mean effects agree between Domainome and Megascale at buried sites (rho 0.84) and not at exposed
sites (rho $-$0.09): exposed substitutions to F/W/Y/I/L/V score worst in the abundance assay, which is consistent with
degron and aggregation effects, whereas in Megascale the worst substitutions are those to P/G/D. Every model with
Megascale-fitted reference energies loses its within-site ranking on Domainome (within-site rho 0.24 for the
model with no reference energies, against $-$0.01 to +0.10 for the models with them).

\textbf{The diagnostic.} Zeroing the reference energies at inference, with the same weights and with no FireProt data
anywhere, gains +0.020 [+0.009, +0.032] on Domainome and +0.017 [+0.001, +0.024] on the flag-kept subset, which
recovers about half of the gain the joint fit achieves there. It costs everywhere else: FireProt 0.464 to 0.443,
S669-strict 0.336 to 0.308, S461-strict 0.552 to 0.519, Megascale natural-strict 0.702 to 0.671. Refitting
without the reference energies recovers nothing on Domainome (0.412), because the other weights absorb them.

\textbf{Reading.} The Megascale reference energies are not globally wrong. They contain a component that is wrong
specifically for an abundance assay at the surface. In the joint fit the Megascale-specific nuisance exceeds
0.4 kcal/mol for one residue only, cysteine ($-$0.55; the next largest is arginine at +0.19), so the assay
contamination is concentrated rather than diffuse.

\textbf{Decision.} Domainome is never trained on and is used for validation only; training on it would make every
Domainome number self-referential.



\section{S16. Test-set looks and selection optimism}

Hyper-parameter selection (term set, transforms, electrostatics variant, cap, loss) used Megascale validation
and grouped CV only, and those tables contain no test column. The model-level decisions were not that clean.
Counting from the version history, \textbf{at least 12 evaluations touched test data}: the first gate evaluation of
the two engine models; the post-audit re-issue of the reporting protocol; audit 1; the AlphaFold treatments; the
residual analysis; the B6 term set; the evolution mode; the final model choice; the elastic-network study; the
differentiable-parameter study; the second critique's phase A; and the shipping decision that audit 2 reversed.
\textbf{At least five changed what ships:} dropping \texttt{site\_scsc} (S10), shipping the joint family, shipping the
AlphaFold variant, and the two reversals after audit 2 (the family choice and the disulfide term).

With roughly 25 model variants and eight benchmarks, the implied optimism on \emph{marginal} deltas is of the same
order as the deltas themselves: the shipped default minus the previous rung is +0.005 [$-$0.004, +0.019], the
AlphaFold variant minus the default is +0.007 [$-$0.003, +0.014], and the joint model minus FoldX on Domainome is
+0.020 [$-$0.002, +0.045]. The headline comparison against FoldX (+0.063) is an order of magnitude larger than any
of these and is the one claim that does not depend on this issue. The validation and CV columns are reported
next to every decision for this reason.



\section{S17. Hand-typed scale tables, replaced}

The first version of the N-cap term used a hand-typed "Doig--Baldwin-type" table from memory. It could not be
verified against the source, and its ranking contradicted the paper's abstract (it placed Gln mid-scale; Doig
and Baldwin 1995 report Gln as the worst N-cap residue). It was replaced by a statistical potential derived from
our own PISCES training split by the method of Richardson and Richardson (1988), E(aa) = $-$ln[P(aa $|$ N-cap) /
P(aa)] in kT units over 58,800 N-caps. The derived ranking S < D < N < T < P < G < H < C < \ldots{} < L < V < I matches
the published ordering, and the model gain is unchanged by the swap (validation 0.7329 to 0.7422 with the term
alone). Two further hand-typed scales (Pace--Scholtz helix propensity, Minor--Kim beta propensity) failed
selection and are unused; both are marked "do not cite" in the code.



\section{S18. The evolution mode: positive, but a different information source}

\textbf{What was done.} One extra weight on a sequence-profile log-ratio, from an HHblits search of UniClust30
(2018\_08) with DDGun's settings; 1,492 chains profiled. No learned model and no embedding, only counting
statistics plus one fitted weight ($-$0.113 $\pm$ 0.007 kcal/mol per nat). Selection on training CV and validation
only.

\textbf{Numbers.} Megascale all 28 per-protein rho 0.685 to 0.726 (+0.040 [$-$0.005, +0.071]); Domainome 0.411 to 0.486
(+0.074 [+0.059, +0.084]); FireProt +0.022 (not significant); S669-strict Spearman +0.037 [+0.004, +0.075];
abundance +0.040. The gain is largest at exposed natural sites, where the structural terms are weakest, and it
is uniform across pLDDT bands on Domainome, where it lifts the sub-70 band from about zero to +0.06.

\textbf{The caveat that keeps it out of the headline.} On the 14 homology-clean natural Megascale proteins the median
change is \textbf{$-$0.006 [$-$0.037, +0.080]}, with a mean per-protein change of +0.014 and 9 of 14 proteins improving;
the seven test proteins that do have a training homologue gain +0.062 (7 of 7 up). One weight cannot memorise
families, so this is either noise or a sign that well-sampled families also have the most informative profiles ---
but the clean subset shows no established gain. In addition, UniClust30 contains homologues of every natural
benchmark protein, which is inherent to any evolutionary method and is not label leakage, but it does mean that
"strict" here only means "no training homologue in the fit".

\textbf{Decision.} Offered as an optional mode, reported separately, never in the headline.



\section{S19. Things we did not try, or could not}

\begin{itemize}
\item \textbf{ColabFold models for viral proteins.} T4 lysozyme, M13 g3p, TEV protease and several others have no
  AlphaFold database entry, so 308 FireProt and 368 Ssym rows are missing from the paired crystal/AlphaFold set.
  Running a structure predictor was outside the scope set for this project.
\item \textbf{The Dunbrack 2010 rotamer library.} Free for academic use but under a manually approved licence that cannot
  be redistributed. We derived our own library instead; tmol bundles a copy whose licence we did not resolve.
\item \textbf{More $\Delta$$\Delta$G data.} ThermoMutDB and ProTherm-derived sets overlap FireProt, S669 and Ssym heavily, so adding
  them would contaminate the benchmarks rather than enlarge the training set.
\item \textbf{Local relaxation.} Never tested. It is the one untried idea with a clear mechanism behind it, since
  Valanciute et al. (2023) show that Rosetta's local Cartesian minimisation is what keeps accuracy on model
  structures where FoldX collapses, and S4 shows random backbone noise is not a substitute.
\item \textbf{Continuous torsion refinement} (gate G4) was never reached, because the dense backbone-dependent library
  plus Boltzmann averaging makes sub-rotamer relaxation largely redundant.
\item \textbf{A deep-learning comparison run by us.} ThermoMPNN, SPURS and RaSP were not rerun; all deep-learning numbers
  in the main text are literature values from other papers' evaluations and are not directly comparable to our
  protocol.
\end{itemize}



\section{S20. Datasets, splits and leakage statistics}

\begin{center}\scriptsize
\begin{tabular}{@{}>{\raggedright\arraybackslash}p{0.107\linewidth}>{\raggedright\arraybackslash}p{0.147\linewidth}>{\raggedright\arraybackslash}p{0.044\linewidth}>{\raggedright\arraybackslash}p{0.198\linewidth}>{\raggedright\arraybackslash}p{0.404\linewidth}@{}}
\toprule
dataset & rows & proteins & use & leakage against our training data \\
\midrule
Megascale (Dataset 3 singles) & 271,231 & 298 & train 215,731 / 239; val 27,328 / 31; test 28,172 / 28 & 9 of 28 test domains have a train homologue (1TUC and 2KXD are 100\% identical to a train domain); strict test subset 19 domains, 14 natural \\
Megascale censored rows & 14,656 & 294 & train only, as intervals (11,363 used) & as above \\
Megascale doubles & 162,472 & 153 & additivity probe only & as above \\
FireProtDB & 3,403 & 99 & train 2,669 / 57; val 387 / 15; test 347 / 27 & 2 of 27 test proteins have a Megascale-train homologue; 3 share a single-linkage family with FireProt trainval (14 rows) \\
S669 & 669 & 94 & external test, strict subset 488 / 87 & 7 of 94 proteins have a Megascale-train homologue (27.1\% of rows); 51 rows homologous to FireProt trainval \\
S461 & 461 & 48 & external test, strict subset 290 / 43 & 113 rows (24.5\%) are mutations that also occur in Megascale-train \\
Ssym & 684 (342 pairs) & 15 & antisymmetry diagnostic only & 96.5\% of rows have a FireProt-trainval homologue; 8 of 15 proteins are exact sequence matches \\
S2648 & 2,633 & 132 & not used in the paper & 68.6\% of rows homologous to FireProt trainval \\
Human Domainome & 536,164 (532,501 with a construct model) & 522 domains & validation / external test & 151 of 517 domains have a Megascale-train homologue, 90 a FireProt-trainval homologue; strict subset 347 domains \\
ProteinGym abundance & 37,947 & 7 assays & external test & 0 of 7 proteins has a homologue in either training set; a median over 7 proteins is not interpretable \\
Paired crystal / AlphaFold & 4,012 usable (3,858 scored) & 198 entities & AlphaFold agreement & 64\% of its rows are FireProt trainval, so it is only an honest test for the Megascale-only models \\
PISCES statistical tables & 9,266 chains & --- & rotamer library, Ramachandran, environment, N-cap & contains homologues of 10/28 Megascale-test, 24/27 FireProt-test, 70/94 S669 and 262/517 Domainome proteins; re-deriving without the 300 overlapping chains moves predictions by 0.0014 kcal/mol and per-protein rho by at most 0.0005 \\
\bottomrule
\end{tabular}\end{center}

Homology is defined by MMseqs2 all-against-all search at 25\% identity with 50\% coverage of the shorter sequence,
taken as single-linkage connected components; 1,228 proteins form 436 families.



\section{S21. Full per-assay results with confidence intervals}

Megascale, Domainome and abundance: per-protein (per-domain) median Spearman rho. FireProt, S669 and S461:
Pearson r. Brackets are 95\% cluster-bootstrap intervals. Rows marked (subset) are computed on fewer rows than
the rest of the column.

\subsection*{Megascale-test, 21 natural proteins (22,159 variants)}

\begin{center}\scriptsize
\begin{tabular}{@{}llll@{}}
\toprule
model & value [95\% CI] & mean per protein & proteins better than FoldX \\
\midrule
PhysDDG AF mode & 0.699 [0.669, 0.727] & 0.705 & 20 / 21 \\
PhysDDG default (Stage 2) & 0.692 [0.659, 0.722] & 0.702 & 19 / 21 \\
PhysDDG joint variant & 0.690 [0.646, 0.718] & 0.691 & 19 / 21 \\
PhysDDG fast (Stage 1) & 0.678 [0.656, 0.711] & 0.691 & 18 / 21 \\
simple baseline B3+ & 0.648 [0.631, 0.670] & 0.654 & 16 / 21 \\
FoldX 5.1 (raw and clipped) & 0.629 [0.551, 0.697] & 0.629 & --- \\
tmol, refitted & 0.622 [0.590, 0.678] & 0.636 & 15 / 21 \\
DDGun3D & 0.516 [0.457, 0.584] & 0.528 & 0 / 21 \\
KORPM & 0.500 [0.483, 0.575] & 0.512 & 3 / 21 \\
\bottomrule
\end{tabular}\end{center}

Other metrics on the same rows: default model Pearson 0.701, RMSE 0.776 kcal/mol, pooled Spearman 0.668;
clipped FoldX 0.596 / 1.845 / 0.592; raw FoldX 0.445 / 3.218 / 0.592; B3+ 0.646 / 0.831 / 0.624; tmol-refit
0.629 / 0.846 / 0.601.

\subsection*{Paired differences, Megascale-test (identical rows)}

\begin{center}\scriptsize
\begin{tabular}{@{}>{\raggedright\arraybackslash}p{0.285\linewidth}>{\raggedright\arraybackslash}p{0.205\linewidth}>{\raggedright\arraybackslash}p{0.205\linewidth}>{\raggedright\arraybackslash}p{0.205\linewidth}@{}}
\toprule
comparison & natural (21) & natural-strict (14) & all (28) \\
\midrule
default $-$ FoldX-clipped & +0.063 [+0.006, +0.120] & +0.062 [+0.009, +0.119] & +0.058 [+0.017, +0.121] \\
AF mode $-$ FoldX-clipped & +0.070 [+0.013, +0.122] & +0.064 [+0.021, +0.119] & +0.065 [+0.023, +0.127] \\
fast $-$ FoldX-clipped & +0.050 [$-$0.001, +0.109] & +0.050 [+0.003, +0.114] & +0.051 [+0.004, +0.113] \\
default $-$ B3+ & +0.044 [+0.013, +0.072] & +0.053 [+0.015, +0.068] & +0.043 [+0.014, +0.066] \\
default $-$ tmol-refit & +0.070 [+0.032, +0.097] & +0.064 [+0.007, +0.118] & +0.066 [+0.015, +0.088] \\
default $-$ base Stage 2 & +0.005 [$-$0.004, +0.019] & +0.004 [$-$0.007, +0.009] & +0.006 [$-$0.004, +0.019] \\
AF mode $-$ base Stage 2 & +0.012 [+0.000, +0.027] & +0.006 [$-$0.001, +0.012] & +0.013 [+0.001, +0.026] \\
fast $-$ base Stage 2 & $-$0.008 [$-$0.011, +0.006] & $-$0.007 [$-$0.027, +0.006] & $-$0.001 [$-$0.019, +0.008] \\
default $-$ FoldX-clipped, Pearson & +0.105 [+0.068, +0.135] & +0.111 [+0.056, +0.148] & +0.097 [+0.064, +0.127] \\
\bottomrule
\end{tabular}\end{center}

\subsection*{The other assays}

\begin{center}\scriptsize
\begin{tabular}{@{}>{\raggedright\arraybackslash}p{0.097\linewidth}>{\raggedright\arraybackslash}p{0.035\linewidth}>{\raggedright\arraybackslash}p{0.054\linewidth}>{\raggedright\arraybackslash}p{0.054\linewidth}>{\raggedright\arraybackslash}p{0.054\linewidth}>{\raggedright\arraybackslash}p{0.092\linewidth}>{\raggedright\arraybackslash}p{0.054\linewidth}>{\raggedright\arraybackslash}p{0.135\linewidth}>{\raggedright\arraybackslash}p{0.265\linewidth}@{}}
\toprule
assay & n / proteins & default & fast & AF mode & joint & B3+ & FoldX & other \\
\midrule
FireProt-test, Pearson & 347 / 27 & 0.451 [0.224, 0.619] & 0.444 [0.208, 0.614] & 0.451 [0.223, 0.620] & 0.478 [0.244, 0.644] & 0.432 [0.182, 0.606] & 0.411 / clipped 0.417 [0.241, 0.568] (subset, 341) & tmol-refit 0.477 (subset, 344); DDGun3D 0.488 [0.223, 0.639]; KORPM 0.397 \\
S669-strict, Pearson & 488 / 87 & 0.336 [0.158, 0.507] & 0.346 [0.163, 0.522] & 0.336 & 0.319 [0.139, 0.479] (470 rows) & 0.339 [0.139, 0.530] & not run & DDGun3D 0.347 [0.171, 0.505]; KORPM 0.194 \\
S461-strict, Pearson & 290 / 43 & 0.552 [0.364, 0.660] & 0.572 [0.366, 0.673] & 0.552 & 0.531 [0.326, 0.647] (285 rows) & 0.570 [0.332, 0.660] & not run & DDGun3D 0.552; KORPM 0.543 \\
Domainome & 509,388 / 493 & 0.418 [0.377, 0.464] & 0.414 [0.376, 0.459] & 0.425 [0.382, 0.468] & 0.450 [0.398, 0.492] & 0.389 [0.352, 0.419] & 0.430 [0.367, 0.483] & KORPM 0.302; DDGun3D 0.282 (subset, 47 domains) \\
Domainome strict & 358,755 / 347 & 0.391 [0.347, 0.433] & 0.389 [0.350, 0.429] & 0.397 [0.350, 0.438] & 0.403 [0.352, 0.468] (320 domains) & 0.369 [0.330, 0.407] & 0.385 [0.310, 0.464] & KORPM 0.274 \\
Domainome, flagged variants withheld & 441,283 / 493 & 0.449 [0.400, 0.481] & 0.450 [0.402, 0.480] & 0.454 [0.403, 0.487] & 0.480 [0.419, 0.510] & 0.407 [0.362, 0.431] & 0.458 [0.396, 0.503] & KORPM 0.315 \\
abundance DMS & 37,947 / 7 & 0.469 [0.343, 0.490] & 0.468 & 0.478 [0.351, 0.499] & 0.484 & 0.446 & not run & a median over 7 proteins is not interpretable; no delta on this assay should be called significant \\
\bottomrule
\end{tabular}\end{center}

Selected paired differences on the other assays: default $-$ FoldX on Domainome $-$0.012 [$-$0.038, +0.021] (strict
+0.007 [$-$0.038, +0.042]); joint $-$ FoldX on Domainome +0.020 [$-$0.002, +0.045]; joint $-$ base Stage 2 on Domainome
+0.038 [+0.019, +0.053]; default $-$ FoldX-clipped on FireProt +0.036 [$-$0.069, +0.102]; default $-$ base Stage 2 on
FireProt $-$0.005 [$-$0.011, $-$0.000]. No FireProt comparison between any pair of methods is significant on any row
set.

\subsection*{Matched rows, where strictness differs between model families}

\begin{center}\scriptsize
\begin{tabular}{@{}>{\raggedright\arraybackslash}p{0.223\linewidth}>{\raggedright\arraybackslash}p{0.202\linewidth}>{\raggedright\arraybackslash}p{0.202\linewidth}>{\raggedright\arraybackslash}p{0.274\linewidth}@{}}
\toprule
 & S669, both-strict (470 / 82) & S461, both-strict (285 / 42) & Domainome, both-strict (323,643 / 320) \\
\midrule
B3+ / base Stage 2 & 0.330 / 0.315 & 0.569 / 0.556 & 0.366 / 0.381 \\
default / fast / AF mode & 0.316 / 0.328 / 0.316 & 0.549 / 0.569 / 0.549 & 0.384 / 0.380 / 0.393 \\
joint, Stage 2 / Stage 1 / AF & 0.319 / 0.330 / 0.319 & 0.531 / 0.550 / 0.531 & 0.403 / 0.405 / 0.410 \\
FoldX-clipped / KORPM / DDGun3D & --- / 0.194 / 0.347 & --- / 0.543 / 0.552 & 0.378 / 0.274 / --- \\
\bottomrule
\end{tabular}\end{center}

\subsection*{Known weaknesses, with numbers}

\begin{center}\scriptsize
\begin{tabular}{@{}>{\raggedright\arraybackslash}p{0.330\linewidth}>{\raggedright\arraybackslash}p{0.150\linewidth}>{\raggedright\arraybackslash}p{0.420\linewidth}@{}}
\toprule
property & PhysDDG & better method \\
\midrule
Ssym antisymmetry, correlation(direct, $-$reverse) / bias & 0.23--0.25 / +1.4 kcal/mol & KORPM 0.88 / +0.15; DDGun3D 0.99 / +0.04 \\
precision@5 for stabilising mutations & 0.15--0.23 & FoldX 0.219; tmol-refit 0.181 (our AUROC is higher: 0.74 against 0.72) \\
buried small-to-large substitutions, per-protein rho & 0.49 & FoldX 0.55 \\
S669-strict, Pearson & 0.336 & DDGun3D 0.347; our own 10-weight baseline 0.339 \\
Domainome below pLDDT 70 & about 0 & every method is about 0 \\
\bottomrule
\end{tabular}\end{center}



\section{S22. Parameter table}

\begin{center}\scriptsize
\begin{tabular}{@{}>{\raggedright\arraybackslash}p{0.089\linewidth}>{\raggedright\arraybackslash}p{0.094\linewidth}>{\raggedright\arraybackslash}p{0.094\linewidth}>{\raggedright\arraybackslash}p{0.047\linewidth}>{\raggedright\arraybackslash}p{0.089\linewidth}>{\raggedright\arraybackslash}p{0.110\linewidth}>{\raggedright\arraybackslash}p{0.318\linewidth}@{}}
\toprule
model & weights & reference energies & intercept & nested base model & total fitted to ddG & also fitted, dropped at inference \\
\midrule
default (Stage 2) & 20 (1 pinned at 0) & 20 & 1 & --- & \textbf{41} (40 effective) & --- \\
fast (Stage 1) & 22 & 20 & 1 & --- & \textbf{43} & --- \\
AF mode & 21 & 20 & 1 & 41 & \textbf{83} & --- \\
joint variants & 19--21 & 20 & 1 & 40 (AF) & 40 / 42 / 81 & 22 (FireProt scale and offset, 20 Megascale reference deltas) \\
\bottomrule
\end{tabular}\end{center}

Fitted weights of the default model (kcal/mol per unit, $\pm$ cluster-bootstrap standard error over 200 draws):

\begin{center}\scriptsize
\begin{tabular}{@{}>{\raggedright\arraybackslash}p{0.138\linewidth}>{\raggedright\arraybackslash}p{0.121\linewidth}>{\raggedright\arraybackslash}p{0.640\linewidth}@{}}
\toprule
term & weight & reading \\
\midrule
\texttt{e\_lj\_atr} & +0.091 $\pm$ 0.005 & Lennard-Jones attraction lost \\
\texttt{e\_lj\_rep} & +0.150 $\pm$ 0.026 & steric clash, signed-log transformed \\
\texttt{e\_solv\_apolar} & +0.107 $\pm$ 0.014 & EEF1 apolar desolvation \\
\texttt{e\_solv\_polar} & +0.129 $\pm$ 0.005 & EEF1 polar desolvation \\
\texttt{e\_hb\_sc\_bb} & +0.223 & side-chain to backbone hydrogen bonds lost \\
\texttt{e\_hb\_sc\_sc} & +0.116 & side-chain to side-chain hydrogen bonds lost \\
\texttt{e\_hb\_bb\_bb} & +0.340 & backbone hydrogen bonds lost \\
\texttt{e\_elec\_dh} & +0.343 $\pm$ 0.045 & Debye--Hückel electrostatics at 0.15 M \\
\texttt{dTS} & $-$0.080 $\pm$ 0.014 & side-chain conformational entropy (negative = stabilising) \\
\texttt{dprior\_avg} & +0.120 $\pm$ 0.020 & rotamer prior \\
\texttt{d\_env\_cn\_ss\_bb} & +0.221 $\pm$ 0.018 & environment statistical potential, per kT \\
\texttt{d\_rama} & +0.205 & Ramachandran at the site \\
\texttt{d\_rama\_prev} & +0.120 & pre-proline map at i$-$1 \\
\texttt{bur} & +0.916 $\pm$ 0.077 & burial, 1 $-$ RSA \\
\texttt{occ\_dvol\_loss} & +1.086 $\pm$ 0.133 & cavity creation \\
\texttt{gly\_posphi} & $-$0.387 $\pm$ 0.082 & Gly$\rightarrow$X at positive phi; a correction to \texttt{d\_rama}, read jointly \\
\texttt{ncap\_pref} & +0.294 $\pm$ 0.037 & N-cap potential, per kT \\
\texttt{helix\_dipole} & +0.086 $\pm$ 0.021 & charge change at a helix terminus \\
\texttt{pro\_nhhb} & 0.000 & at its bound; the backbone hydrogen-bond term already carries the lost N--H \\
\texttt{ss\_broken} & 0.000 & pinned; Megascale contains no disulfide-breaking rows \\
\bottomrule
\end{tabular}\end{center}

The AlphaFold model adds \texttt{scale\_plddt} = $-$0.966 (bootstrap mean $-$1.12, 2.5--97.5 percentile $-$2.52 to $-$0.23 over
200 draws), which scales the prediction by 1 $-$ 0.97$\cdot$(1 $-$ pLDDT/100) and is identically zero on crystal input.

Removing every sign constraint moves all other weights by at most 0.022 kcal/mol and changes validation and CV
metrics by at most 0.0001.

Tables estimated from structures and never from ddG: the backbone-dependent rotamer library (PISCES, 9,266
chains), Ramachandran maps, the environment potential (about 180 entries), the N-cap potential (20 entries), and
about 15 frozen CHARMM19 and EEF1 constants with literature provenance. The repulsive radii carry one structural
calibration factor of 0.9, chosen by native-rotamer recognition and not by ddG.



\section{S23. FoldX protocol}

\begin{itemize}
\item \textbf{Binary:} FoldX 5.1 under an academic licence, run outside this repository, protocol RepairPDB followed by
  BuildModel on the benchmark structures.
\item \textbf{Sign:} FoldX reports positive = destabilising, which matches our convention; no flip was applied.
\item \textbf{Options not recorded.} \texttt{numberOfRuns}, \texttt{--temperature}, \texttt{--pH}, \texttt{--ionStrength} and \texttt{--vdwDesign} were not
  written down and no configuration files were kept. The binary's defaults are \texttt{numberOfRuns 1}, 298 K, pH 7,
  0.05 M, \texttt{vdwDesign 2}, and that is what we assume. Every statement about the FoldX baseline should therefore
  read "FoldX 5.1, RepairPDB then BuildModel, run options not recorded, defaults assumed". With a single
  BuildModel run some of FoldX's error is run-to-run variance that averaging would remove, which makes our
  comparison mildly pessimistic for FoldX on rank.
\item \textbf{pH.} FireProt rows span pH 1.5 to 11, so a single pH 7 run is wrong for part of that set --- for FoldX and
  for us equally, since our model has no pH term.
\item \textbf{Structures.} Megascale has no crystal structures: all 298 domains are the AlphaFold2 models released with
  the assay, so FoldX ran RepairPDB on AlphaFold2 models and our model is fitted on the same files. Both methods
  see exactly the same input, so the comparison is fair in the identical-rows sense; it is the generalisation
  reading that needs the caveat.
\item \textbf{Coverage.} FoldX covers 100\% of Megascale-test, 341 of 347 FireProt-test rows and 493 of 518 Domainome
  constructs. We use its construct-level Domainome run, not the full-length one, because the assay measures
  isolated constructs.
\item \textbf{Transform symmetry.} We report FoldX raw and clipped to [$-$3, 8], a fixed monotone transform with nothing
  fitted, because its raw predictions reach 71 kcal/mol and that alone costs 0.15 Pearson. Our own repulsion
  feature carries a signed-log transform that was selected on validation. Compressing the extreme values helps both models, so
  the clipped column is the like-for-like comparison.
\end{itemize}



\section{S24. Audit findings, summarised}

Two audits were run by agents that did not write the code that they audited. Both recomputed metrics from the raw
prediction files with their own code. Their reports are released unedited.

\textbf{Audit 1} (leakage, evaluation correctness, over-claims):

\par\noindent 1. The catalytic-CXXC exception was derived from two FireProt \emph{test} rows (S11). Removing it drops the joint
   models' FireProt Pearson from 0.48 to 0.39.
\par\noindent 2. S461 is contaminated: 113 of 461 rows (24.5\%) are mutations that also occur in Megascale-train. Report strict
   subsets only.
\par\noindent 3. Ssym is not a test set for any model trained on FireProt (85.7\% exact duplicate mutations).
\par\noindent 4. The pre-registered gate against FoldX fails on rank for the two early engine models (+0.048 [$-$0.008, +0.105]
   and +0.058 [$-$0.003, +0.113] on the primary subset) and passes only on the two secondary subsets. The honest
   wording is "ahead in point estimate on every Megascale subset, significantly so on two of three".
\par\noindent 5. The large Pearson advantage over FoldX is mostly its uncapped tail: clipping lifts FoldX from 0.445 to 0.596,
   so the gap is +0.10 and not +0.25.
\par\noindent 6. The headline natural subset is not homology-clean (21 proteins, 14 strict), although the strict number is not
   lower.
\par\noindent 7. The statistical-table overlap with benchmark proteins is real but negligible (per-protein rho changes by at
   most 0.0005 when the tables are re-derived without the 300 overlapping chains).
\par\noindent 8. Main-table cells were not on identical rows; the paired comparisons were.
\par\noindent 9. The wild-type-to-wild-type test was tautological and was replaced by tests of the property that matters.

\textbf{Audit 2} (the joint models, inference correctness, claims), with every model-card cell recomputed
independently:

\par\noindent 1. \textbf{Critical.} The FireProt column was computed on 338 of 347 rows and the two missing rows were exactly the
   ones that discriminate the models (S11). Fixed: all FireProt numbers are now on 347 rows.
\par\noindent 2. \textbf{Critical.} The shipped model had been selected partly on test numbers; at least 12 test evaluations and at
   least 5 that changed what ships (S16). Fixed: the family was re-decided on validation, CV and transfer only,
   and the 12 evaluations are disclosed.
\par\noindent 3. \textbf{Critical.} The \texttt{final\_*} prediction columns in the repository were stale and belonged to a superseded
   family. Fixed: all shipped files are now produced by one export script and the prediction columns were
   regenerated.
\par\noindent 4. Domainome had no homology filter; every absolute number falls by about 0.03 when training homologues are
   removed, and the abundance assays had never been screened (they turn out to be clean, 0 of 7).
\par\noindent 5. The parameter count was understated: 41 / 43 / 83, not 40.
\par\noindent 6. The guarantee that crystal input is never attenuated holds on every shipped path, but it rests on a string
   match with a B-factor fallback that would misclassify 21 of 94 S669, 9 of 48 S461 and 5 of 99 FireProt
   crystal structures if a caller used the lower-level entry point without declaring the structure type.
\par\noindent 7. Confidence intervals quoted to three decimals from 200 bootstrap draws; the headline reproduces at 2,000
   draws (+0.063 [+0.008, +0.096] for the joint model against clipped FoldX), but marginal bounds should be
   rounded to two decimals.
\par\noindent 8. The paired crystal/AlphaFold set is 64\% joint-model training data, so the "AlphaFold is never worse than the
   crystal" result holds only for the Megascale-only models, and even there the intervals for the fitted models
   touch zero.
\par\noindent 9. A list of claims the final report must not make, which we have followed: no "best structure-only model on
   FireProt", no identical-rows claim without a fixed strictness rule, no homology-clean Domainome or abundance
   claim, no "the physics generalises across assays", no "one test look", no "40 parameters", no significance
   for the marginal deltas listed in S16, and no AlphaFold-specific correction that improves accuracy.

Inference-side checks that passed: the nuisance parameters are never applied at inference (verified by
corrupting them and observing bit-identical output); the command-line tool reproduces the stored predictions to
7e$-$6; a wild-type-to-wild-type substitution returns exactly zero; per-term contributions sum to the prediction
to machine precision; predictions are bitwise reproducible across runs and across GPUs; signs are correct and
well separated (32 buried hydrophobic-to-charged variants of one protein average +2.76 kcal/mol, all positive).

\end{document}
